\documentclass[12pt,a4paper]{article}
\usepackage[utf8]{inputenc}
\usepackage[T2A]{fontenc}
\usepackage[russian]{babel}
\usepackage{amsmath,amssymb,amsthm}
\usepackage{booktabs,array}
\usepackage{geometry}
\usepackage{microtype}
\geometry{margin=2.3cm}

\newtheorem{protocol}{Протокол}
\newtheorem{proposition}{Предложение}
\newtheorem{nogocriterion}{Критерий провала}

\title{Проверка гипотезы арифметического магнита $\pi$}
\author{}
\date{4 августа 2026 г.}

\begin{document}
\maketitle

\section{Идея проверки}

Ранее отношение масс тау-лептона и мюона сравнивалось с формулой
\begin{equation}
 \frac{m_\tau}{m_\mu}
 =\pi^2+2\pi+\frac23-\frac{\alpha}{3}.
\end{equation}
В замороженном списке из $1485$ коротких выражений эта формула заняла первое
место. Однако первое место в списке ещё не означает малую вероятность
случайного совпадения. Если числовая сетка плотная, у каждого числа найдётся
свой победитель.

Проверяется грамматика
\begin{equation}
 f(x;n,r,c)=x^2+nx+r+c\alpha,
\end{equation}
где $n=0,\ldots,4$, $r=p/q$ при $|p|\leq4$, $q\leq6$, а
$c\in\{0,\pm1,\pm\tfrac12,\pm\tfrac13,\pm\tfrac23\}$. Сначала полагается
$x=\pi$, затем $x$ свободно пробегает интервал $[2,4]$.

\begin{protocol}[Три контроля]
Используются три независимых вопроса:
\begin{enumerate}
 \item какую долю числового интервала покрывают окрестности всех $1485$ формул;
 \item какая доля оснований $x\in[2,4]$ даёт совпадение не хуже, чем $\pi$;
 \item исчезает ли эффект после запрета формул сложнее заявленной.
\end{enumerate}
Ширина попадания фиксируется заранее равной фактической ошибке
$\delta=0.0006679209$.
\end{protocol}

\section{Результат}

Точное объединение интервалов, без статистического моделирования, даёт:
\begin{center}
\begin{tabular}{>{\raggedright\arraybackslash}p{0.48\textwidth}r}
\toprule
Контроль & Доля попаданий \\
\midrule
Полный диапазон значений грамматики & $9.25\%$ \\
Окно ширины $2$ вокруг отношения $m_\tau/m_\mu$ & $13.02\%$ \\
Полный диапазон, сложность не выше заявленной & $5.86\%$ \\
Основания $x\in[2,4]$, совпадение не хуже $\pi$ & $9.08\%$ \\
То же при сложности не выше заявленной & $5.40\%$ \\
\bottomrule
\end{tabular}
\end{center}

Таким образом, результат $\pi$ находится примерно в лучших девяти процентах
оснований, но не в исключительной области. Даже после ограничения сложности
примерно одно основание из девятнадцати справляется не хуже.

Есть и наглядный контроль:
\begin{align}
 e^2+2e+4-\alpha
 &=16.8183224,\\
 (\sqrt{10})^2+2\sqrt{10}+\frac12-\alpha
 &=16.8172580.
\end{align}
Обе формулы в той же грамматике ближе к измеренному отношению, чем формула с
$\pi$. Это не делает $e$ или $\sqrt{10}$ физическими объяснениями. Наоборот,
это показывает, насколько легко грамматика производит убедительные числа.

\begin{proposition}[Арифметический магнит]
Внутри замороженной постфактум-грамматики утверждение ``формула с $\pi$
занимает первое место'' не является сильным статистическим свидетельством.
Вероятность столь же точного попадания имеет порядок нескольких процентов, а
не долей промилле.
\end{proposition}

\section{Что именно произошло}

Без члена $c\alpha$ полный диапазон покрыт лишь на $1.07\%$. После добавления
девяти малых сдвигов покрытие возрастает до $9.25\%$. Значит, главным
``магнитом'' является не только $\pi$, а связка
\begin{equation}
 \text{крупный каркас }x^2+nx+r
 \quad+\quad
 \text{тонкая ручка настройки }c\alpha.
\end{equation}
Малые рациональные коэффициенты при $\alpha$ работают как зубья частой
числовой расчёски.

\begin{nogocriterion}[Честный вердикт]
Проверка не доказывает, что формула ложна. Она доказывает более узкое
утверждение: одна численная близость не отличает физический закон от удачного
арифметического выбора. Для возврата физического статуса формула должна быть
получена из заранее заданного оператора и затем успешно предсказать новое
число, которое не участвовало в выборе грамматики.
\end{nogocriterion}

Все числа воспроизводятся программой
\texttt{s2t\_pi\_arithmetic\_magnet\_audit.py}; полный протокол сохранён в
\texttt{s2t\_pi\_arithmetic\_magnet\_results.json}.

\end{document}